\section{Related Work}
\label{sec:related}

\paragraph{Agentic Tool-use benchmarks.}
Several benchmarks evaluate LLM tool use. ToolQA~\citep{zhuang2023toolqa} tests question answering over external data sources; API-Bank~\citep{li2023api} evaluates API selection across hundreds of real APIs; Toolformer~\citep{schick2023toolformer} and ToolLLM~\citep{qin2023toolllm} train models to invoke tools correctly; Gorilla~\citep{patil2024gorilla} targets correct API call generation; and BFCL~\citep{patil2025berkeley} provides a comprehensive function-calling leaderboard spanning single-turn, multi-turn, and agentic settings. All these benchmarks evaluate \emph{whether} models can use tools correctly, assuming every task requires a tool. In contrast, \benchmark{} evaluates the tool-call \emph{decision}: given the correct tool, the model must decide whether to use it or solve directly (Table~\ref{tab:benchmark_compare}).

\paragraph{Efficient tool calling.}
Recent work on reducing unnecessary tool calls takes several approaches. \citet{xu2025alignment} fine-tune the model to invoke tools only when its confidence is low, reducing calls by $\sim$50\% on arithmetic and QA tasks. \citet{wu2025joint} jointly optimize agent instructions and tool descriptions via verbalized feedback, reducing calls by up to 70\%. \citet{yang2026toward} survey efficiency across memory, tool learning, and planning in LLM agents. These works directly build costly interventions, such as SFT pipelines or iterative prompt optimization, without first investigating \emph{why} models overcall. We instead take a mechanistic approach in a more realistic agentic setting, showing that the model's hidden state already encodes the tool-call decision and can be leveraged with only a lightweight linear probe that takes seconds to train and adds negligible inference cost.

\paragraph{Probing and controlling LLM behavior.}
Linear probing has revealed that LLM hidden states encode syntactic structure~\citep{hewitt2019structural}, factual knowledge~\citep{burns2022discovering}, truthfulness~\citep{li2023inference}, and self-knowledge~\citep{kadavath2022language}. Building on these findings, a growing body of work uses internal representations to steer model behavior: activation addition~\citep{turner2024activation} and representation engineering~\citep{zou2023representation} modify hidden activations during generation, concept bottleneck LLMs~\citep{sun2024concept} route predictions through interpretable concept layers, and recent work edits model weights guided by steering directions for reasoning control~\citep{sun2025thinkedit, sun2026steer2edit, yan2025reflctrl} and skill unlearning~\citep{li2025effective}. Our work extends probing to a new domain, agentic tool-use decisions, showing that models ``mostly know'' when they need tools but fail to act on it. Unlike prior steering methods that modify activations or weights, we steer via output prefilling based on a probe prediction, requiring no modification to the forward pass and remaining compatible with any serving infrastructure.
